% Artigo de teste do fluxo LaTeX -> HTML -> Medium.
% Usa as mesmas macros do template Tá no Hype (README do repositório).

\hypearticle
{Do LaTeX ao Medium: testando um fluxo de publicação}
{Equipe Hype Data \& AI}
{

\hspace*{1.5em} Este artigo existe para validar uma ideia simples: escrever uma única vez em
\LaTeX{} e distribuir em dois formatos, o PDF da newsletter e uma página web que o Medium
consiga importar mantendo título, imagens, listas e destaques.

\subtitle{Por que uma fonte única?}

Quando o mesmo texto vive em dois lugares, cedo ou tarde eles divergem. Com uma fonte única,
a correção de uma vírgula no \texttt{.tex} se propaga para o PDF e para a versão web na
próxima compilação. Acentuação, símbolos como 100\% e \& também precisam sobreviver à
conversão.

\articleimage{0.8}{img/figura.png}{Figura 1. Etapas do fluxo de publicação testado neste artigo.}

\subtitle{O que queremos verificar}

\begin{hypeitemize}
  \item O título e o subtítulo chegam ao Medium com a hierarquia correta;
  \item As imagens são baixadas pelo importador (URLs absolutas e públicas);
  \item Listas, citações e caixas de destaque mantêm o sentido;
  \item O link canônico aponta para a página de origem.
\end{hypeitemize}

\highlight{Se o Medium importar este bloco como uma citação, o destaque roxo do PDF
virou um \emph{pull quote} na versão web.}

\begin{infobox}
\textbf{Caixa informativa.} No PDF ela é laranja; na web, uma citação com classe própria.
Um exemplo de fórmula em linha: $E = mc^2$.
\end{infobox}

\subtitle{Próximo passo}

Com a página publicada, o último passo é manual: importar a URL pelo recurso
\emph{Import a story} do Medium e conferir o resultado com a lista de verificação do
repositório.

\originalsource{https://github.com/hype-usp/ta-no-hype-newsletter}

}
